\section{Object-Centric and Structured Physical Dynamics}
\label{sec:object_centric_physical_dynamics}

Object-centric physical dynamics respond to a central limitation of pixel-level prediction: physical interaction is often easier to describe in terms of entities, relations, and forces than in terms of all pixels at once. The previous section showed that action-conditioned and stochastic video predictors can roll images forward, but their image-space objectives do not by themselves expose object identity, pairwise interaction, hidden causes, or planning-relevant state. Structured dynamics models introduce an explicit or learned intermediate representation in which objects carry state and relations mediate change. This shift does not remove the need for perception or generation; rather, it changes the predictive question from "what should the next frame look like?" to "which entities interact, how do their states change, and how can those changes be rendered, evaluated, or used for decisions?"

\subsection{Objects and Relations as Physical Abstractions}

The basic object-centric premise is that a physical scene can be modeled as a set of entities whose future states depend on their own attributes, external effects, and interactions with other entities. Interaction Networks make this premise explicit by representing a system as an attributed graph and decomposing dynamics into relation-centric interaction functions and object-centric update functions \cite{battaglia2016_interaction_networks_prov}. This graph view supplies two useful inductive biases for world modeling. First, it separates object state from pairwise interaction, which makes the learned dynamics more interpretable than a monolithic pixel predictor. Second, it supports compositional generalization: the completed notes report that Interaction Networks generalize across different numbers and configurations of objects in controlled physical systems \cite{battaglia2016_interaction_networks_prov}.

This abstraction is powerful but also restrictive. Interaction Networks assume explicit object states and graph structure rather than discovering objects directly from video, so they are better viewed as a dynamics template than as a complete visual world model. The same limitation clarifies why object-centric video models are difficult: a useful structured simulator must solve perception, object binding, relational dynamics, and often rendering or planning at the same time. This is the point at which structured physical dynamics connects back to video prediction. Pixel-level models avoid object discovery but struggle to expose stable physical state; graph-based models expose physical structure but often assume the structure is already available.

\subsection{From Visual Perception to Relational Dynamics}

Visual Interaction Networks bridge this gap by connecting raw visual observations to relational physical prediction. In the completed notes, VIN combines a CNN-based visual encoder with an interaction-network dynamics predictor, producing object-factorized latent state codes that are rolled forward by relational dynamics \cite{watters2017_visual_interaction_networks_prov}. This architecture is important for the survey because it does not force a choice between pixels and structured state: perception maps images into object-centric codes, and dynamics operates over those codes. The result is a visually grounded model of physical interaction rather than a purely symbolic simulator.

The same design also shows why structured dynamics can do more than forward extrapolation. VIN is reported to infer hidden physical causes such as invisible objects and unknown mass from visual effects, and to outperform pixel-to-state and privileged state-to-state baselines on long rollout position error in its controlled settings \cite{watters2017_visual_interaction_networks_prov}. These findings support a cautious claim: object-centric visual dynamics can help expose latent physical variables when the environment has a stable entity structure. However, the evidence remains bounded by synthetic or controlled scenes with supervised state targets and limited visual complexity. Scaling this style of model to cluttered, partially observed, and open-ended real video remains an open problem rather than an established result.

\subsection{Object-Oriented Prediction, Rendering, and Planning}

Object-oriented prediction becomes more world-model-like when the learned representation is used not only for rollout, but also for planning. O2P2 combines perception, pairwise physics, and rendering around object-centric latent vectors, and trains the system through image reconstruction and prediction losses rather than direct labels for object properties \cite{janner2019_o2p2_prov}. This design connects structured physical dynamics to generative simulation: the model must encode objects from visual input, predict an interaction outcome in representation space, and render the predicted configuration back into pixels.

The planning result in O2P2 sharpens the evaluation lesson from video prediction. The completed notes report that O2P2's learned object-oriented representation supports downstream planning better than an object-agnostic video predictor baseline in a tower-building task \cite{janner2019_o2p2_prov}. The important point is not merely that object-centric models reconstruct scenes, but that their abstractions can be useful for choosing actions in structured physical tasks. At the same time, O2P2 should not be overstated as a general simulator: it uses segmentation masks or object proposals, predicts a steady-state outcome rather than a full temporal rollout, and is evaluated mainly in structured block-stacking settings.

\subsection{Slot-Based and Compositional Extensions}

Several later or adjacent lines are natural extensions of this section, but the current evidence base does not yet support detailed technical claims about them. The paper index includes contrastive structured world models, slot-based visual dynamics, entity abstraction for visual model-based reinforcement learning, visual predictive models for billiards, compositional object-based dynamics, and unsupervised learning of object structure and dynamics from videos. These categories are directly relevant to object-centric and physical dynamics because they address latent object discovery, slot dynamics, compositional simulation, or video-based structure learning. However, completed notes are still needed before this section can state what each method empirically demonstrates. Claims about C-SWM, SlotFormer, OP3, billiards-style visual physics, compositional object dynamics, or unsupervised object-structure learning should therefore remain [NEEDS SOURCE].

The evidence-ready papers nevertheless identify the comparison dimensions that these extensions should eventually fill in. A structured model may assume ground-truth object states, infer object-centric codes from pixels, use segmentation masks or proposals, learn slots without direct object labels, or induce keypoints and latent entities from video. These choices affect compositionality, interpretability, and scaling. Explicit graphs are interpretable but require object definitions; visual encoders connect to pixels but can depend on supervised state targets; object-oriented renderers close the loop back to image space but may rely on proposals or simple geometries. A mature account of slot-based and compositional video dynamics will need to compare these assumptions directly once the missing notes are complete.

\subsection{Strengths, Limits, and Bridge to Latent Control}

The main strength of object-centric physical dynamics is that it makes interaction structure explicit. Compared with pixel-level prediction, object and relation representations can express which entities interact, factor dynamics across pairs or groups, expose hidden physical variables in controlled settings, and provide abstractions that are more directly useful for planning \cite{battaglia2016_interaction_networks_prov,watters2017_visual_interaction_networks_prov,janner2019_o2p2_prov}. These strengths explain why the object-centric branch is central to a survey on predictive world models from video: it offers a principled route from visual prediction to physical abstraction.

The limitations are equally important for the survey's argument. Current evidence-ready object-centric models often rely on synthetic or controlled environments, explicit object states, supervised state targets, segmentation masks, object proposals, or restricted interaction settings \cite{battaglia2016_interaction_networks_prov,watters2017_visual_interaction_networks_prov,janner2019_o2p2_prov}. Even when object-like structure emerges from robot video prediction, the structure can remain weakly object-centric and image-space-bound \cite{finn2016_physical_interaction_video_prediction_prov}. The broader challenge is to learn object structure, hidden variables, and interaction dynamics with less supervision while preserving usefulness for planning or control \cite{finn2016_physical_interaction_video_prediction_prov,battaglia2016_interaction_networks_prov,watters2017_visual_interaction_networks_prov,janner2019_o2p2_prov,bruce2024_genie_prov}.

This limitation motivates the next section on latent world models for planning and control. Object-centric models make physical structure explicit, but they can be hard to learn robustly from complex real videos. Latent world models take a different route: they learn compact predictive states that may not correspond to human-interpretable objects, but can support planning or policy learning efficiently. The transition is therefore not a replacement of structured dynamics by latent dynamics, but a shift in emphasis from explicit physical decomposition to task-usable predictive state.

% Concise Claim-Evidence Audit:
% - Major claims:
%   Object-centric dynamics use entities and relations as inductive biases for physical prediction; visual interaction models bridge pixels and relational dynamics; object-oriented representations can support planning in structured tasks; object discovery and scaling remain open problems.
% - Supporting papers:
%   P007 / battaglia2016_interaction_networks_prov; P008 / watters2017_visual_interaction_networks_prov; P009 / janner2019_o2p2_prov; limited bridge evidence from P005 / finn2016_physical_interaction_video_prediction_prov and P013 / bruce2024_genie_prov.
% - Weak or missing evidence:
%   C-SWM, SlotFormer, OP3, billiards visual physics, compositional object dynamics, and unsupervised object-structure-from-video claims remain [NEEDS SOURCE] until notes for P017--P019 and P029--P031 are completed.
% - Revision TODOs:
%   Add a comparison table after missing notes are complete; verify whether slot-based and compositional models should receive their own subsection; revisit limitations after broader realistic-video evidence is available.
